\section{Numerical simulation}
\label{sec:performance-evaluation}

This section presents a performance evaluation of magic state cultivation with lattice surgery (MSC-LS).
We use Stim~\cite{gidney2021stim} and PyMatching~\cite{higgott2025sparse} for the evaluation.
Our implementation is available at \hiddenforreview{\url{https://github.com/yutakahirano/msc-ls/}}.
Because Stim can only simulate Clifford circuits, we replace $T$ gates in the injection and cultivation stages with $S$ gates.
This modification causes the distillation circuit to output $S\ket{+}$ instead of $T\ket{+}$.
Gidney~\textit{et al.}~\cite{Gidney2024MagicStateCultivation} conjecture that cultivating $T\ket{+}$ has a logical error probability approximately twice as high as cultivating $S\ket{+}$.
We note that, in the parameter regime considered below (\autoref{subsec:performance-evaluation-simulation-settings}), namely $d_\mathrm{color} = 3$ and $p_\mathrm{phys} = 10^{-3}$, the validity of this approximation for the injection and cultivation stages is supported by statevector simulations reported in Ref.~\cite{Gidney2024MagicStateCultivation}.
This approximation may also quantitatively affect complementary-gap-based postselection, and thus the final logical error probability and success probability.
Nevertheless, we adopt it in this section because Clifford simulation with this approximation is standard practice in MSC-style analyses.
Recently, several techniques~\cite{wan2025simulate, surti2025efficient, li2025soft} have been proposed to enable simulations of MSC-like protocols without this approximation.
Applying such techniques to MSC-LS is an important direction for future work.

\subsection{Simulation settings} \label{subsec:performance-evaluation-simulation-settings}
We evaluate the performance of MSC-LS using several parameters.
These include the color code distance $d_{\mathrm{color}} = 3$, the intermediate rotated surface code distance $d_\mathrm{intermediate}$, the final rotated surface code distance $d_\mathrm{final}$, and the number of rounds of syndrome extraction $r$ performed while waiting for the complementary gap computation.
Simulations are carried out under uniform depolarizing circuit noise of strength $p_\mathrm{phys}$.
Unless otherwise specified, we set $d_\mathrm{final} = 11$, $r = 10$, and $p_\mathrm{phys} = 10^{-3}$.
For the original magic state cultivation, we use a grafted code distance of $d_\mathrm{grafted} = 15$, following the observation of Gidney~\textit{et al.}~\cite{Gidney2024MagicStateCultivation} that the logical error probability of the distance-15 grafted code is comparable to that of the distance-11 rotated surface code.

We take a ``round'' as the unit of distillation time.
A round is a series of circuit layers that corresponds to one round of syndrome extraction in the 2D color code or the rotated surface code.
This unit allows us to compare the spacetime overhead of distillation with magic state cultivation.
\autoref{alg:performance-evaluation-spacetime-cost} presents the algorithm for calculating the spacetime overhead of distillation.

\begin{figure}[!t]
  \captionsetup{type=algorithmfigure}
  \hrule \vspace{10pt}
  \caption{Calculate the spacetime overhead}
  \label{alg:performance-evaluation-spacetime-cost}
  \hrule \vspace{2pt}
  \begin{algorithmic}
    \State $cost \gets 0$
    \For{each round}
      \State $q \gets$ the number of qubits used in this round
      \State $p \gets$ the success probability of this round
      \State $cost \gets (cost + q) / p$ 
    \EndFor
    \State \textbf{return} $cost$
  \end{algorithmic}
  \vspace{2pt} \hrule
\end{figure}

\subsection{Results}
\label{subsec:performance-evaluation-results}

\begin{figure}[t!]
\centering
\includegraphics[width=8cm]{fig/performance-evaluation-error-detection.pdf}
\caption{Logical $X$ and $Z$ error probabilities introduced by the lattice surgery procedure.}
\label{fig:performance-evaluation-error-detection}
\Description{Fully described in the text.}
\end{figure}

\autoref{fig:performance-evaluation-error-detection} illustrates $X$ and $Z$ logical error probabilities resulting from the lattice surgery procedure described in \autoref{subsec:msc-ls-ls} as a function of $p_\mathrm{phys}$.
Here, to focus on the logical errors caused by lattice surgery, the simulation runs the following sequence.
First, an ideal $\ket{+}_L$ ($\ket{0}_L$) state is encoded in the distance-3 2D color code.
Next, the lattice surgery procedure described in \autoref{subsec:msc-ls-ls} is performed to teleport the quantum state to the rotated surface code of distance $d_\mathrm{intermediate} = 5$.
Finally, we apply one round of noise-free syndrome extraction on the rotated surface code, followed by postselection and a noise-free logical measurement in the $X$ ($Z$) basis, after which we verify the outcome.
This evaluates the $Z$ ($X$) logical error probability of the procedure.
Fitting the data to the form $C\,p_{\mathrm{phys}}^{\alpha}$ gives $\alpha \approx 2.92$ for $Z$ errors and $\alpha \approx 3.30$ for $X$ errors, indicating that the procedure preserves a fault distance of 3, consistent with expectations.

\begin{figure}[t!]
\centering
\includegraphics[width=8.5cm]{fig/performance-evaluation-end-to-end.pdf}
\caption{
Logical error probability of the resulting magic state as a function of distillation success probability.
Blue: MSC-LS with $d_\mathrm{intermediate} = 3$; orange: MSC-LS with $d_\mathrm{intermediate} = 5$; green: MSC.
Here we do not use the lookup table for MSC-LS.
}
\label{fig:performance-evaluation-end-to-end}
\Description{Plot of logical error probability versus distillation success probability for three protocols.
The horizontal axis shows distillation success probability, and the vertical axis shows logical error probability on a logarithmic scale.
Three increasing curves are shown.
The blue and green curves lie close to each other over much of the plotted range, indicating similar logical error probabilities.
The orange curve lies below them over much of the range, indicating lower logical error probability at the same distillation success probability.
For a fixed logical error probability, the orange curve extends farther to the right than the others, indicating higher distillation success probability.}
\end{figure}

\autoref{fig:performance-evaluation-end-to-end} illustrates the trade-off between the logical error probability and end-to-end distillation success probability of MSC-LS and MSC.
Each point in a series corresponds to a different complementary gap threshold used in the final postselection.
MSC-LS achieves logical error probabilities comparable to those of MSC but outperforms it in end-to-end distillation success probability.
Moreover, MSC-LS with $d_\mathrm{intermediate} = 5$ outperforms MSC-LS with $d_\mathrm{intermediate} = 3$.
% We attribute the small residual difference in the minimum achievable logical error probability between MSC and MSC-LS with $d_\mathrm{intermediate} = 5$ primarily to skipping syndrome extraction for the six-weight stabilizer $X_{0145ab}$ in the lattice surgery procedure (\autoref{subsec:msc-ls-ls}).

The logical error probability in the injection and cultivation stages is $2 \times 10^{-7}$, compared to $8 \times 10^{-7}$ for the lattice surgery procedure (shown in \autoref{fig:performance-evaluation-error-detection}) and $1.6 \times 10^{-6}$ for the end-to-end protocol (shown in \autoref{fig:performance-evaluation-end-to-end}).
This is consistent with the observation in Ref.~\cite{Gidney2024MagicStateCultivation} that the escape stage dominates the logical error budget.

To evaluate the sensitivity of these results to the noise model, we vary each error component independently while keeping the others fixed.
\autoref{fig:performance-evaluation-sensitivity-more} and \autoref{fig:performance-evaluation-sensitivity-less} show the resulting logical error probability as a function of the distillation success probability when the error probability of each component is increased and decreased, respectively.
These results show how each error source affects MSC-LS separately, providing a more detailed view of which hardware improvements would most directly improve the performance of the protocol.
In particular, the logical error probability is more sensitive to changes in the idling and two-qubit gate error probabilities than to changes in the other error probabilities.

\begin{figure}[t!]
\centering
\includegraphics[width=8.5cm]{fig/performance-evaluation-sensitivity-more.pdf}
\caption{
Logical error probability of the resulting magic state as a function of distillation success probability.
The orange curve represents the baseline noise model, which is uniform depolarizing circuit noise of strength $10^{-3}$.
The other curves represent modified noise models in which one error component is set to $1.5 \times 10^{-3}$ while the other error probabilities are kept at their baseline values.
}
\label{fig:performance-evaluation-sensitivity-more}
\Description{Plot of logical error probability versus distillation success probability for the baseline noise model and modified noise models with increased single-qubit gate, two-qubit gate, reset, measurement, and idling error probabilities.
The horizontal axis shows distillation success probability, and the vertical axis shows logical error probability on a logarithmic scale.}
\end{figure}

\begin{figure}[t!]
\centering
\includegraphics[width=8.5cm]{fig/performance-evaluation-sensitivity-less.pdf}
\caption{
Logical error probability of the resulting magic state as a function of distillation success probability.
The orange curve represents the baseline noise model, which is uniform depolarizing circuit noise of strength $10^{-3}$.
The other curves represent modified noise models in which one error component is set to $6.7 \times 10^{-4}$ while the other error probabilities are kept at their baseline values.
}
\label{fig:performance-evaluation-sensitivity-less}
\Description{Plot of logical error probability versus distillation success probability for the baseline noise model and modified noise models with decreased single-qubit gate, two-qubit gate, reset, measurement, and idling error probabilities.
The horizontal axis shows distillation success probability, and the vertical axis shows logical error probability on a logarithmic scale.}
\end{figure}


\begin{figure}[t!]
\centering
\includegraphics[width=8.5cm]{fig/performance-evaluation-lookup-table-coverage.pdf}
\caption{
Coverage of the lookup table.
The horizontal axis denotes the ratio of distillation attempts that would be additionally rejected by the lookup table.
The vertical axis represents the number of entries in the table.
}
\label{fig:performance-evaluation-lookup-table-coverage}
\Description{Plot of lookup table size versus coverage.
The horizontal axis shows coverage, and the vertical axis shows the number of entries in the lookup table.
The curve rises gradually at low coverage and then increases sharply near coverage 0.096, reaching about 160,000 entries.}
\end{figure}


\begin{figure}[t!]
\centering
\includegraphics[width=8.5cm]{fig/performance-evaluation-distillation-lifetime-without-lookup-table.pdf}
\caption{
Qubit usage and survival probability during distillation without the lookup table.
Each column represents a round.
The blue region represents the survival probability of distillation.
The orange region represents the physical qubit usage normalized to the maximum qubit usage (241 qubits in this case). 
}
\label{fig:performance-evaluation-distillation-lifetime-without-lookup-table}
\Description{Fully described in the text.}
\end{figure}

\begin{figure}[t!]
\centering
\includegraphics[width=8.5cm]{fig/performance-evaluation-distillation-lifetime-with-lookup-table.pdf}
\caption{
Qubit usage and survival probability during distillation with the lookup table, using the same color scheme as \autoref{fig:performance-evaluation-distillation-lifetime-without-lookup-table}.
}
\label{fig:performance-evaluation-distillation-lifetime-with-lookup-table}
\Description{Combined step plot of qubit usage and survival probability during distillation with the lookup table.
Because additional lookup-table-based postselection is performed before code expansion, the survival probability is lower from that point until the final postselection.
The final survival probability is unchanged.}
\end{figure}

We construct lookup tables using \autoref{alg:construct-lookup-table} with various complementary gap thresholds $g_\mathrm{th}$.
For each $g_\mathrm{th}$ setting, we sample $10^{10}$ shots, with the threshold for invalid cases set to $c_\mathrm{th} = 100$.
\autoref{fig:performance-evaluation-lookup-table-coverage} shows the coverage of the lookup table with $g_\mathrm{th} = 15.48$.
For this table, the offline construction took about 5.7 hours using 90 CPU cores on Intel Xeon Platinum 9242 processors, corresponding to approximately 510 core-hours.
At maximum size, the table contains 159,321 entries and rejects 9.6\% of distillation attempts.
\autoref{fig:performance-evaluation-distillation-lifetime-without-lookup-table} and \autoref{fig:performance-evaluation-distillation-lifetime-with-lookup-table} show the qubit usage and survival probability during distillation.
Without the lookup table (\autoref{fig:performance-evaluation-distillation-lifetime-without-lookup-table}), 43\% of samples are rejected before code expansion, 22\% by the final postselection, and 35\% are accepted.
By contrast, with the lookup table (\autoref{fig:performance-evaluation-distillation-lifetime-with-lookup-table}), 52\% of samples are rejected before code expansion, 13\% are rejected by the final postselection, and 35\% are accepted.
This early rejection reduces the spacetime overhead by 15\% in this setting.

% Our lookup-table construction is conservative because it uses a large number of Monte Carlo samples and retains only syndrome patterns observed at least $c_\mathrm{th}$ times.
% We quantify this conservative nature using the false-rejection rate, defined as the fraction of all distillation attempts that are rejected by the lookup table but would have passed the final postselection.
% For example, for $g_\mathrm{th} = 15.48$, a Monte Carlo simulation with $2 \times 10^9$ shots yielded 706,469,472 accepted shots with the correct logical output, 1,153 accepted shots with a logical error, and 1,293,529,375 rejected shots when only the final postselection was applied.
% When the lookup-table-based early rejection was also applied, the same simulation yielded 706,469,365 accepted shots with the correct logical output, 1,153 accepted shots with a logical error, and 1,293,529,482 rejected shots.
% Thus, there were 107 false rejections, corresponding to an observed false-rejection rate of $107/(2 \times 10^9) = 5.35 \times 10^{-8}$, which is negligibly small.

\begin{figure}[t!]
\centering
\includegraphics[width=8.5cm]{fig/performance-evaluation-distillation-overhead.pdf}
\caption{
Spacetime overhead of MSC-LS with and without the lookup table and MSC.
}
\label{fig:performance-evaluation-distillation-overhead}
\Description{Plot of logical error probability versus spacetime overhead of distillation for MSC-LS without table, MSC-LS with table, and MSC.
The horizontal axis shows spacetime overhead in qubit-rounds, and the vertical axis shows logical error probability on a logarithmic scale.
All three curves decrease as spacetime overhead increases.
The two MSC-LS curves lie to the left of the MSC curve, and the lookup-table version lies slightly to the left of the version without the table in the lower logical-error region.}
\end{figure}

\autoref{fig:performance-evaluation-distillation-overhead} depicts the spacetime overhead of distillation of MSC-LS and MSC.
It shows that in the logical error probability range of $[10^{-6}, 10^{-5}]$, MSC-LS requires less than half the spacetime overhead required for MSC.
For example, MSC-LS reduces the spacetime overhead by 59\% (55\%) when distilling an $S\ket{+}$ state with a logical error probability of $2 \times 10^{-6}$ ($4 \times 10^{-6}$).
The lookup table performs more effectively in the lower range of logical error probability, which is expected, given that postselection is performed more strongly in this range.

\autoref{fig:performance-evaluation-lookup-table-coverage} also shows that the lookup table is useful even when the number of entries in the table is limited.
For example, with only 5,000 entries, the coverage reaches 8\%.
Hence, the lookup table is useful even in memory-constrained systems.
The number of detectors contained in a table entry is 54 with our simulation setting, and hence each table entry requires only 54 bits and is very lightweight.